\section{Introduction}

A vision-language model (VLM) can answer an image question correctly and yet
fail to \emph{update} that answer when a single, decision-relevant visual factor
changes. Standard accuracy on a static image conflates two abilities: producing a
correct answer for the observed scene, and revising the decision when the scene
is intervened upon. We argue that the second ability---\emph{decision consistency
under controlled interventions}---is the property that matters for trustworthy
perception, and that it is under-measured by accuracy on fixed images.

CertVIC measures decision consistency under controlled, single-factor,
real-image interventions. From a real image and an object mask we synthesize a
minimally edited counterpart that changes the expected answer (e.g.\ removing the
object beneath a supported item) or, for control items, changes nothing
decision-relevant. Neutral prompts query the model on the original and edited
images; an item is \emph{consistent} when the original/edited answer pair
respects the intended change. We define the
\emph{intervention-consistency gap} $\Delta = a - p$, where $a$ is original-image
accuracy and $p$ is the consistency rate, and we \emph{certify} a lower bound on
$\Delta$ using an anytime-valid confidence sequence (CS), so that the guarantee
holds even under sequential, optional-stopping evaluation on a small budget.

Two design commitments distinguish CertVIC. First, it is \emph{budget-aware and
zero-cost}: the entire pipeline runs on local CPU or free GPU with open-source
models and public data, and it certifies results with a stated statistical
guarantee rather than relying on large fixed-$n$ test sets. Second, it is
\emph{recipe-first}: we release source pointers, hashes, masks, edit
parameters, and scripts rather than redistributing pixels, so results are
reproducible without licensing hazards.

\paragraph{Contributions.}
(1) A controlled, single-factor real-image intervention pipeline with explicit
edit-validity and human review gates.
(2) An anytime-valid certification of the intervention-consistency gap, with
power planning and optional-stopping diagnostics.
(3) Construct-validity baselines and ablations (text-only, caption-only,
single-image, control edits) that show the task cannot be solved without seeing
the change.
(4) A zero-cost, recipe-first, reproducible open-model evaluation workflow.
Empirical values remain [RESULT REQUIRED] until eligible pilot and main runs
exist.
